\documentclass{article}

\usepackage{amsthm}
\usepackage{graphicx}
% \usepackage{amssymb}

\newtheorem{proposition}{Proposition}

\title{MATH290 Problem 1}
\author{Caleb Hessing}
\date{September 6, 2026}

\begin{document}
\maketitle 

A famous pirate captain has hidden his treasure in three chests. One chest contains
gold coins, one chest contains counterfeit coins, and one chest contains a mixture of
gold and counterfeit coins. Each chest is labeled to indicate its contents: either G
for gold, C for counterfeit, or B for both. Because this pirate captain is very clever
and dishonest, none of the labels match the actual contents of the chest. You cannot
visually distinguish a gold coin from a counterfeit one, but you can perform a lengthy
chemical test to tell the two apart. What is the minimum number of coins you need to
test to determine the contents of all three chests, and what strategy should you use?

\begin{proposition}
    Only one coin test is necessary to determine the contents of all three chests.
\end{proposition}

\begin{proof}\nobreakspace

    Notice that because the contents of the chests do not match their labels,
    the chest labelled B must contain either all gold coins or all
    counterfeit coins. Therefore, we can determine chest B's contents
    by testing just one of its coins. If the one coin is counterfeit, all of chest
    B chest's coins are also; if it is gold, all of chest B's coins are gold.
    
    Consider the case where the coin we test is counterfeit, revealing that
    chest B contains only counterfeit coins. That means that of two remaining chests, labelled 
    C and G, one contains only gold coins and the other contains both
    gold and counterfeit coins. Because the chests' contents don't match their labels,
    chest G cannot contain only gold; therefore, it must contain both counterfeit and gold 
    coins. That leaves only one chest unidentified---chest C---which, having
    eliminated all other options, must contain only gold coins.

    Consider the opposite case, where the coin we test reveals chest B
    contains only gold coins. Following the same logic, we can deduce that
    chest C must contain both counterfeit and gold coins, and chest G must 
    contain only counterfeit coins.

    In either case, it takes one coin to determine the contents of chest B,
    and zero coins to deduce the contents of chests C and G. Therefore,
    only one coin test is necessary to determine the contents of all three chests.

\end{proof}

\end{document}